\section{Introduction}\label{sec:introduction}

All graphs in this paper are finite, simple, and undirected. A
\emph{correspondence assignment} on a graph $G$ consists of pairwise
disjoint finite sets $L(v)$, one for each vertex, and a matching
$M_{uv}\subseteq L(u)\times L(v)$ for each edge $uv$. A coloring of the
assignment chooses $c(v)\in L(v)$ at every vertex so that
$(c(u),c(v))\notin M_{uv}$ for every edge $uv$. The matchings need not
agree with one another around cycles. The \emph{correspondence chromatic
number}, also called the \emph{DP-chromatic number}, $\chi_{\mathrm{DP}}(G)$,
is the least integer $k$ for which every correspondence assignment with
$|L(v)|\ge k$ admits a coloring.

This notion includes list coloring: replace each color in each list by
a distinct copy, and match copies of the same color along an edge.
Ordinary coloring is obtained by giving every vertex the same list
before taking these copies. Thus
\[
 \chi(G)\le \chi_{\mathrm{list}}(G)\le\chi_{\mathrm{DP}}(G)
 \le \Delta(G)+1,
\]
where the last inequality follows by greedy coloring. We prove a
logarithmic improvement when the clique number is bounded.

\begin{theorem}\label{thm:main}
For every integer $r\ge4$ there are constants $C_r>0$ and $\Delta_r\ge3$
such that every $K_r$-free graph $G$ of maximum degree
$\Delta\ge\Delta_r$ satisfies
\[
 \chi_{\mathrm{DP}}(G)
 \le \left\lceil C_r\frac{\Delta}{\log\Delta}\right\rceil.
\]
The constants depend only on $r$.
\end{theorem}

In particular, the estimate is uniform in the number of vertices,
the sizes of the lists, and the edge matchings. Throughout the paper,
all logarithms are natural. We do not optimize the constant or the
degree threshold.

For a fixed finite graph $F$, call $G$ \emph{$F$-free} if it contains no
subgraph isomorphic to $F$; the subgraph need not be induced.

\begin{corollary}\label{cor:forbidden}
For every finite graph $F$ there are constants $C_F>0$ and
$\Delta_F\ge3$ such that every $F$-free graph $G$ of maximum degree
$\Delta\ge\Delta_F$ satisfies
\[
 \chi_{\mathrm{DP}}(G)
 \le \left\lceil C_F\frac{\Delta}{\log\Delta}\right\rceil.
\]
\end{corollary}
\begin{proof}
Set $r=\max\{4,|V(F)|\}$. Since $K_r$ contains $F$ as a subgraph,
every $F$-free graph is $K_r$-free. Apply Theorem~\ref{thm:main}.
\end{proof}

The ordinary-coloring version of Corollary~\ref{cor:forbidden} is
Conjecture~3.1 of Alon, Krivelevich, and Sudakov~\cite{AKS1999}.
The corollary therefore proves that conjecture in the stronger
correspondence-coloring form. If $F$ contains a cycle, the order
$\Delta/\log\Delta$ cannot be improved in general: graphs with girth
larger than $|V(F)|$ can have maximum degree at most $\Delta$ and
chromatic number of this order~\cite{Bollobas1978}; see also the proof of
\cite[Proposition~1.2]{AKS1999}.

\subsection{Historical context}
List coloring was introduced independently by Vizing~\cite{Vizing1976}
and by Erd\H{o}s, Rubin, and Taylor~\cite{ErdosRubinTaylor1980}.
Dvo\v{r}\'ak and Postle introduced correspondence coloring in their
2015 preprint, published as~\cite{DvorakPostle}, to allow the
identification of colors to vary from edge to edge. This flexibility
was central to their application to list coloring planar graphs.
An arbitrary correspondence assignment need not admit globally
consistent color names, so a list-coloring argument that uses
independence across such names requires additional work.

An early step toward logarithmic coloring bounds imposed a stronger
restriction on short cycles. Kim~\cite{Kim1995} proved
$\chi_{\mathrm{list}}(G)\le(1+o(1))\Delta/\log\Delta$ for graphs
of girth at least five. Johansson subsequently obtained the list-coloring bounds
$O(\Delta/\log\Delta)$ for triangle-free graphs and
$O_r(\Delta\log\log\Delta/\log\Delta)$ for $K_r$-free graphs.
His triangle-free argument is presented by Molloy and
Reed~\cite[Chapter~13]{MolloyReed2002}; the extended version of
Bansal, Gupta, and Guruganesh~\cite[Sections~6--9]{BansalGuptaGuruganesh2018}
develops his recursive probability modifications and entropy estimates.
Molloy~\cite[Theorems~1 and~2]{Molloy2019} obtained the
triangle-free list bound with leading constant $1+o(1)$ and gave a new
proof of the clique-free bound, with explicit dependence on $r$.

Bernshteyn~\cite{Bernshteyn2016} extended Johansson's triangle-free
bound to correspondence coloring, and later extended Molloy's
asymptotic triangle-free bound and the clique-free bound to that
setting~\cite[Theorems~1.3 and~1.5]{Bernshteyn2019}. In particular,
the previous general upper bound for fixed $r\ge4$ was
\[
 \chi_{\mathrm{DP}}(G)
 =O_r\!\left(\frac{\Delta\log\log\Delta}{\log\Delta}\right).
\]
Davies, Kang, Pirot, and Sereni~\cite[Theorem~30]{DKPSLocalOccupancy}
developed a framework based on local occupancy in the hard-core model
and improved the leading constants. Przyby\l o~\cite[Theorem~10]{Przybylo2024}
also obtained bounds under clique exclusion in the graph whose vertices
are list entries and whose edges are forbidden pairs, while retaining
the maximum degree of the underlying graph as the degree parameter.
These results retain the
$\log\log\Delta$ factor for a general fixed forbidden clique.

Other families of forbidden subgraphs admit the sharper order.
Alon, Krivelevich, and Sudakov~\cite{AKS1999} proved that a graph
whose neighborhoods each span at most $\Delta^2/f$ edges, with
$2\le f\le\Delta^2$, has chromatic number $O(\Delta/\log f)$.
Vu~\cite{Vu2002} extended this local-sparsity bound to list coloring.
Alon, Krivelevich, and Sudakov also deduced
$O_F(\Delta/\log\Delta)$ for every fixed almost bipartite forbidden
graph $F$, meaning that deleting at most one vertex from $F$ leaves a
bipartite graph. Anderson, Bernshteyn, and
Dhawan~\cite{AndersonBernshteynDhawan2023,AndersonBernshteynDhawan}
proved correspondence-coloring bounds with leading constants
$1+o(1)$ for bipartite forbidden graphs and $4+o(1)$ for almost
bipartite forbidden graphs. Their more general theorems assume that
the graph of forbidden pairs excludes $F$ and use its maximum degree
in place of $\Delta(G)$.

More recently, Dhawan, Janzer, and
Methuku~\cite[Theorems~1.4 and~1.5]{DhawanJanzerMethuku} proved
$O_t(\Delta/\log\Delta)$ for ordinary coloring of $K_{t,t,t}$-free
graphs, and hence for every fixed three-colorable forbidden graph.
They also obtained the independence bound
$(1-o(1))n\log d/d$ in terms of the average degree $d$.
Their coloring theorem
concerns ordinary coloring; it does not supply a list or correspondence
bound. Theorem~\ref{thm:main} removes the $\log\log\Delta$ factor
for correspondence coloring with any fixed forbidden clique $K_r$, $r\ge4$.

The structural part of our argument is related to the independence
problem. Ajtai, Erd\H{o}s, Koml\'os, and
Szemer\'edi~\cite{AEKS1981} asked for a lower bound of order
$n\log d/d$ for $K_r$-free graphs of average degree $d$.
Shearer~\cite{Shearer1983,Shearer1995} obtained the asymptotic
triangle-free bound and a clique-free bound with a
$\log\log d$ loss. The companion manuscript
\emph{A logarithmic independence bound for clique-free graphs}
\cite{OpenAIIndependence} proves the logarithmic bound and develops
the weighted multipliers, neighborhood walks, and vertex splitting
used here. The recursive use of clique-free neighborhoods has a
precedent in~\cite[Section~3]{AEKS1981}; entropy increments controlling
distances between walk distributions have a precedent
in~\cite[Section~2, Theorem~8]{BDKY2015}. We adapt these ingredients
to a local adjustment statement that also bounds how many entries
from each list an operation affects. We establish the required
weighted and dual estimates directly.

The remaining probabilistic ingredients also have established
antecedents. We use the simultaneous coupling bound presented by
Angel and Spinka~\cite{AngelSpinka}, with its finite-state proof
included. Concentration for sums in which each independent input affects
only a bounded number of summands is the setting of
Gavinsky, Lovett, Saks, and Srinivasan~\cite{GavinskyReadK}.
We derive the weighted form needed here from Finner's generalized
H\"older inequality~\cite{Finner} and Hoeffding's exponential-moment
bound~\cite{Hoeffding1963}. The Lov\'asz local
lemma~\cite{ErdosLovasz1975} then makes the local round estimates
hold simultaneously.


\subsection{The proof and its main ingredients}

The proof constructs a partial coloring in small random rounds. At an
uncolored vertex $v$, each available entry $i\in L(v)$ carries a
nonnegative weight $w_i$. Its total weight
$W_v=\sum_{i\in L(v)}w_i$ must remain bounded below, while the individual
weights remain small. These two conditions guarantee many available
entries. The round colors enough neighbors to make the remaining degree
decrease faster than this supply of entries.

Let $H$ be the graph whose vertices are the positive-weight entries and
whose edges are the forbidden pairs in the matchings. The lists partition
$V(H)$ into independent sets. Every vertex of $H$ has at most one
neighbor in each part, and every clique of $H$ projects to a clique in
$G$. Thus $H$ is $K_r$-free. This graph allows us to work directly with
the correspondences, without identifying entries from different lists.

A round independently activates each entry with probability proportional
to its weight. An activated entry can be chosen if none of its neighbors
in $H$ activates. Choosing it removes all conflicting entries from the
remaining lists. The potential controlling the loss of weight has three
terms: an entropy relative to the initial weights, a linear penalty for
lost total weight, and the incident edge mass
\[
 D_v=\sum_{i\in L(v)}\sum_{j\in N_H(i)}w_iw_j.
\]
The last term measures the amount of conflicting weight around $v$.
Deleting newly colored neighbors decreases it and pays for the cost of
the round.

\paragraph{Triangles and local adjustments.}
An activation at a common neighbor of adjacent entries $i,j$ removes
both weights at once. Even after compensating their individual losses
to first order, this shared deletion creates a positive edge
correlation. Summed at an entry $i$, its cost is proportional to
\[
 t_i=w_i\sum_{j\in N_H(i)}w_j
                 \sum_{z\in N_H(i)\cap N_H(j)}w_z.
\]
This is twice the weighted mass of triangles through $i$.
Lemma~\ref{lem:adjustment} supplies random local multiplications and
deletions of weights whose combined contribution offsets $t_i$, up to
a constant times the incident edge mass, simultaneously at all entries.
The lemma also bounds the total rate of operations at each entry.

The structural argument develops the multiplier and entropy-splitting
methods of the companion manuscript
\emph{A logarithmic independence bound for clique-free graphs}
\cite[Sections~2--5]{OpenAIIndependence}. We prove all the ingredients
needed here. Finite-dimensional separation tests the vertex inequalities
and the rate bounds against arbitrary nonnegative prices. This reduces
the adjustment problem to bounding the resulting priced triangle mass
by a constant multiple of the priced edge mass, together with the
allowance supplied by the rate bounds. Because vertex prices may vary
widely, we group triangles by the dyadic size of their median price and
rescale the weights within each level. The dual constraints then bound
the weight of edges between nearby sets. Walks inside neighborhoods
guide a splitting of vertices into copies, preserving most triangle
weight while sharply reducing every neighborhood weight. Repetition
leaves bounded neighborhood weights, where the triangle-to-edge
comparison follows directly. The accumulated losses are controlled
independently of the number of price levels.

\paragraph{Controlling participation within each list.}
For correspondence coloring, a support of small graph radius may still
contain many entries from one list. Radius alone therefore does not
give the concentration needed for the coloring round. The adjustment
lemma controls both quantities: each operation is supported in a ball,
and it affects only a bounded number of entries in each part of the
partition. When the neighborhood weight is $O(\log\Delta)$, both its
multiplier bound and its participation bound are $\Delta^{o(1)}$.

The additional bound comes from the local sets used in the walk
argument. A vertex is added to such a set only when it has a prescribed
positive amount of neighbor weight in the preceding set. The
one-neighbor-per-part property converts this admission threshold into
a bound on the number of vertices added from each part. It survives
splitting because the projection of the split graph is injective on
directed edges. No lower bound on individual vertex weights is needed.
This gives the adjustment lemma as a separate statement for partitioned
weighted graphs, before its application to coloring.

\paragraph{Concentration and completion.}
The trial uses independent flags for activations and adjustments.
Different entry weights after the trial need not be independent.
Instead, Lemma~\ref{lem:read-concentration} controls sums of functions
of independent inputs when each input is used by few summands. The
participation bound applies directly to the total mass and the
non-edge part of the potential at a list. We control the edge part by
concentrating each matching sum separately.

Activations may collide across forbidden pairs. To show that these
collisions cost little, we first bound the degree of their activated
subgraph. Too many collision edges would then contain a large matching,
whose endpoint activations are independent. This supplies the required
tail estimate without a decomposition into globally named colors.
The Lov\'asz local lemma gives simultaneous progress at all vertices.
Trimming extreme weights restores the assumptions for the next round.
Once every remaining degree is smaller than the number of available
entries, a greedy coloring completes the assignment.

Section~\ref{sec:adjustment} states the adjustment lemma and develops
the dual local estimates. Section~\ref{sec:triangles} proves the
entropy and splitting estimates and completes that lemma.
Section~\ref{sec:rounds} constructs the random update, proves its
expectation and concentration bounds, and restores the weight conditions
after simultaneous progress. Section~\ref{sec:completion} iterates this
round and completes Theorem~\ref{thm:main}.

\subsection{Conventions}
Neighborhoods are open; graph-distance balls include their centers.
An edge or triangle has weight equal to the product of its vertex
weights. Unless indicated otherwise, sums over edges and triangles are
unordered. A graph homomorphism is \emph{injective on directed edges}
if distinct ordered adjacent pairs have distinct images. It may
identify vertices that do not occur in the same directed edge.
Constants implicit in $O_r(\cdot)$ and sufficiently large thresholds
may depend on $r$, but not on the graph, lists, or matchings. Auxiliary
graphs may be empty.
